\begin{frame}{``Attention Is Not Explanation'' --- and the Reply}

    {\scriptsize An \textbf{attention} layer assigns a weight to every input token, and those weights
    are routinely displayed as ``what the model looked at.''}

    \begin{columns}[T]
        \begin{column}{0.48\textwidth}
            \begin{tcolorbox}[colback=raiRed!5, colframe=raiRed!70,
                              title=Jain \& Wallace (NAACL 2019)]
                \scriptsize
                Across several NLP tasks:
                \begin{itemize}
                    \item Attention weights are frequently \textbf{uncorrelated} with
                          gradient-based importance.
                    \item One can construct \textbf{very different} attention distributions that
                          give \textbf{the same} prediction.
                \end{itemize}
                Conclusion: standard attention modules do not provide meaningful explanations.
            \end{tcolorbox}
        \end{column}
        \begin{column}{0.48\textwidth}
            \begin{tcolorbox}[colback=raiBlue!5, colframe=raiBlue!60,
                              title=Wiegreffe \& Pinter (EMNLP 2019)]
                \scriptsize
                The conclusion depends on what you mean by ``explanation.'' They propose four
                tests --- including a uniform-weight baseline and a diagnostic on frozen
                weights --- and find that adversarial attention distributions, where they exist,
                do poorly on the diagnostic.

                \vspace{0.3em}
                Attention is not \emph{no} evidence; it is not \emph{sufficient} evidence.
            \end{tcolorbox}
        \end{column}
    \end{columns}

    \begin{block}{What to do in practice}
        \centering
        \scriptsize Use attention maps as a \textbf{debugging view}, never as a justification you
        hand to a user or a reviewer. For a transformer attribution, use a method built for it (Integrated
        Gradients, attention rollout, gradient\,$\times$\,input) and validate it.
    \end{block}

    \vspace{0.1em}
    {\scriptsize S.~Jain and B.~C.~Wallace, ``Attention is not Explanation,'' NAACL 2019,
    \href{https://arxiv.org/abs/1902.10186v3}{arXiv:1902.10186v3}. S.~Wiegreffe and Y.~Pinter,
    ``Attention is not not Explanation,'' EMNLP 2019,
    \href{https://arxiv.org/abs/1908.04626v2}{arXiv:1908.04626v2}.}
\end{frame}
